\section{Introduction}
\label{sec:introduction}

% 第 1 段：问题动机 — 设计流编排是工程师最耗时的工作
Circuit design is less about writing individual components and more about \emph{orchestrating design flows}: a chip tape-out requires traversing schematic capture, SPICE simulation, layout generation, DRC/LVS physical verification, parasitic extraction, timing signoff, and dozens of tool-to-tool handoffs. Each EDA tool is a closed system; where tools refuse to interoperate, the only recourse is hand-written glue---a Tcl driver, a SKILL wrapper, a Python script parsing a simulation log---precisely the intent-to-code synthesis at which LLMs excel. This reframes EDA workflow automation as \emph{agentic flow orchestration}: a coding agent that can call live EDA tools, inspect their outputs, and iteratively refine a deliverable across multiple steps.
% 电路设计的核心不在于编写单个组件，而在于"编排设计流"：一次芯片 tape-out 需要穿越原理图录入、SPICE 仿真、版图生成、DRC/LVS 物理验证、寄生参数提取、时序签核，以及数十次工具间衔接。每件 EDA 工具都是封闭系统；当工具之间无法互操作时，唯一的办法就是手写粘合代码——一段 Tcl 驱动、一个 SKILL 包装、一个解析仿真日志的 Python 脚本——而这正是 LLM 所擅长的"意图到代码"综合。这把 EDA 工作流自动化重新定义为"智能体式流程编排"：一个能够调用真实 EDA 工具、检视其输出、并跨多步迭代改进交付件的编程智能体。

% 第 2 段：现有 benchmark 的缺口
Existing LLM benchmarks for EDA fall into two categories, neither of which captures this orchestration challenge. \emph{Code-generation benchmarks} (RTLLM~\cite{lu2024rtllm}, VerilogEval~\cite{liu2023verilogeval}) evaluate single-function RTL generation from natural language---one module, one testbench, no tool interaction. \emph{Flow-repair benchmarks} (FlowFixer-EDA~\cite{dong2026flowfixer}, EDA-FlowBench) evaluate diagnosing and patching failures in open-source RTL-to-GDS runs---but only on open-source tools (Yosys, OpenROAD), not the commercial EDA ecosystem where most engineers work. Neither evaluates the \emph{end-to-end agentic loop}: read a task description, call live EDA tools through an MCP bridge, inspect execution results, and produce a complete deliverable package. Software engineering benchmarks like SWE-bench~\cite{jimenez2024swebench} test this loop for GitHub issues, but no equivalent exists for EDA design flows.
% 现有的面向 EDA 的 LLM 基准可分为两类，但都未捕捉到这一编排挑战。"代码生成类基准"（RTLLM、VerilogEval）评估从自然语言到单函数 RTL 的生成——一个模块、一个测试台、无工具交互。"流程修复类基准"（FlowFixer-EDA、EDA-FlowBench）评估在开源 RTL-to-GDS 流程中诊断并修复失败——但仅限于开源工具（Yosys、OpenROAD），而非大多数工程师实际使用的商业 EDA 生态。两者都没有评估"端到端智能体回路"：读取任务描述、通过 MCP 网桥调用真实 EDA 工具、检视执行结果、并产出完整的交付件包。软件工程类基准如 SWE-bench 针对 GitHub issue 测试了这一回路，但在 EDA 设计流领域尚无对应物。

% 第 3 段：我们的方法 — 从真实用户轨迹构建 benchmark
We present \bench, built from an unprecedented data source: \textbf{real user interaction traces} of 268 engineers using the Cline~\cite{cline2026} coding agent across 42 workstations in a production IC design environment. Each task in \bench{} is derived from an actual multi-turn session where an engineer used an LLM agent to solve a real design-flow problem---ranging from generating a SPICE deck for a sense amplifier, to writing a DRC rule-check script for a DRAM layout, to automating netlist connectivity tracing for a mixed-signal block. We systematically mined 698 hosts, identified 1{,}511 candidate sessions (match-count $\geq 30$ EDA-keyword hits), and curated \Ntasks{} tasks across 13 EDA domains, each with $\geq 8$ executable assertions and a reference deliverable.
% 我们提出 \bench，其数据来源前所未有：来自生产 IC 设计环境中 42 台工作站上 268 位工程师使用 Cline 编程智能体的真实交互轨迹。\bench{} 中的每个任务都源自一段真实的多轮会话——工程师曾用 LLM 智能体解决一个真实的设计流问题，从为灵敏放大器生成 SPICE 仿真卡，到为 DRAM 版图编写 DRC 规则检查脚本，再到为混合信号模块自动化网表连通性追踪。我们系统性地扫描了 698 台主机，筛出 1,511 个候选会话（EDA 关键词命中数 $\geq 30$），并在 13 个 EDA 领域中精心整理了 \Ntasks{} 个任务，每个任务配有不少于 8 条可执行断言和一份参考交付件。

% 第 4 段：评测协议 + 防泄题
Each task follows a unified \mode{} contract: the agent receives a task brief and input assets, must orchestrate remote EDA tools (on a dedicated EDA server) through an MCP~\cite{lee2025autoeda} bridge, and produce a deliverable package (\texttt{submit.zip}). Scoring is execution-based---a judge unpacks the deliverable and runs reference assertions against it---with a \texttt{PreToolUse} hook that blocks the agent from accessing the answer key (check.py, reference/), ensuring the pass rate reflects genuine agentic capability rather than answer leakage.
% 每个任务遵循统一的 \mode{} 契约：智能体接收任务简述与输入附件，必须通过 MCP 网桥调用远程 EDA 工具（位于专用 EDA 服务器），并产出交付件包（\texttt{submit.zip}）。判分基于执行——判分器解包交付件并运行参考断言——同时配合 \texttt{PreToolUse} 钩子，阻止智能体访问答案（check.py、reference/），确保通过率反映真实的智能体能力而非答案泄露。

% 第 5 段：核心发现
We evaluate a single frontier LLM agent (DeepSeek-V4-Pro) under a controlled Cline harness and obtain a striking result: it achieves only \textbf{10.0\% Pass@1} (62/620 decided runs) on the frozen \Neval{}-run evaluation, far below the 76.5\% reported on single-script EDA benchmarks~\cite{zhu2026edascript}. The gap is domain-dependent: standardized-flow domains where the deliverable is largely a script or tcl template (pnr\_signoff 20.8\%, characterization 19.2\%, eda\_other 18.2\%) pass at roughly twice the overall rate, while domains that require deep live-tool interaction---physical verification (DRC/LVS runs, 2.9\%) and Aether layout editing (2.9\%)---pass at the bottom of the range. We hypothesise that this ordering is explained by \emph{tool-chain grounding depth}---the number of live EDA tool interactions needed to produce a correct deliverable---and we separate what is measured (the domain ordering, the outcome mix, and the per-run trace mechanics) from what remains an untested hypothesis (that grounding depth dominates model scale).
% 我们在受控的 Cline 框架下评测单一前沿 LLM 智能体（DeepSeek-V4-Pro），得到一个引人注目的结果：它在冻结的 \Neval{}-run 评测上仅取得 \textbf{10.0\% Pass@1}（62/620 已判定 run），远低于单脚本 EDA 基准上报告的 76.5\%。
% 该差距随领域而变：交付件主要为脚本或 tcl 模板的标准化流程领域（pnr_signoff 20.8\%、characterization 19.2\%、eda_other 18.2\%）的通过率约为整体的两倍，而需要深度实时工具交互的领域——物理验证（DRC/LVS 运行，2.9\%）与 Aether 版图编辑（2.9\%）——通过率垫底。
% 我们假设这一排序可由"工具链接地深度"——产出正确交付件所需的真实 EDA 工具交互次数——解释；并把"已测量"（领域排序、结果构成、逐 run 轨迹机制）与"仍属未检验假设"（接地深度压过模型规模）明确区分开。

% 第 6 段：贡献
\noindent\textbf{Contributions.}
% \noindent\textbf{贡献。}
\begin{itemize}
    \item \textbf{Benchmark.} \bench: \Ntasks{} tasks from real user traces, 13 EDA domains, execution-based scoring with anti-leakage protocol (Section~\ref{sec:benchmark}).
    % \textbf{基准。} \bench：来自真实用户轨迹的 \Ntasks{} 个任务，覆盖 13 个 EDA 领域，执行式判分加防泄题协议（第~\ref{sec:benchmark}~节）。
    \item \textbf{Evaluation.} First systematic evaluation of an LLM agent on multi-step EDA design-flow orchestration with live commercial tools, revealing a 10.0\% Pass@1 ceiling (62/620 decided runs) on the frozen \Neval{}-run snapshot and a $7.1\times$ point-estimate spread between the best and weakest domain (pnr\_signoff 20.8\% vs.\ physical\_verify/layout\_aether 2.9\%) (Section~\ref{sec:experiments}).
    % \textbf{评测。} 首次系统性地评估 LLM 智能体在真实商业工具上的多步 EDA 设计流编排能力，在冻结的 \Neval{}-run 快照上揭示 10.0\% 的通过率天花板（62/620 已判定 run），以及最好/最弱领域间 7.1 倍的点估计跨度（pnr\_signoff 20.8\% vs.\ physical\_verify/layout\_aether 2.9\%）（第~\ref{sec:experiments}~节）。
    \item \textbf{Analysis.} A failure taxonomy with run-level anchors (PASS 62 / FAIL 558 / ERROR 70 / ABSENT 6), an explicit protocol for per-run attribution, and a trace-anatomy study of 47 archived runs showing where agent effort actually goes (median 26.5 tool calls, 41.2\% through the MCP bridge) (Section~\ref{sec:analysis}).
    % \textbf{分析。} 带 run 级锚点的失败分类法（PASS 62 / FAIL 558 / ERROR 70 / ABSENT 6）、逐 run 归因的显式协议，以及对 47 份归档轨迹的解剖研究，说明智能体的投入实际流向何处（中位 26.5 次工具调用，41.2\% 走 MCP 网桥）（第~\ref{sec:analysis}~节）。
\end{itemize}
